% RecSys Odyssey repository-backed lecture note.
% Add figures with: \coursefigure{figures/name.svg}{Caption}
\section{Signals are evidence, not truth}

\begin{abstract}
A click is evidence of interest, but it is also shaped by what was exposed and where.
\end{abstract}

\subsection{Concept}
Explicit feedback is deliberately supplied: ratings, likes, dislikes or survey answers. Implicit feedback is inferred from behavior: impressions, clicks, watch time, skips, purchases and returns. Stronger events are not automatically cleaner. A purchase can reflect price or availability; a long watch can reflect autoplay. Always log the impression as well as the response, otherwise “not clicked” is confused with “never shown.”

\begin{equation*}
interaction = (user, item, event, timestamp, context, exposure)
\end{equation*}

\subsection{Key terms}
\begin{itemize}
\item \textbf{Explicit feedback.} A preference the user intentionally reports.
\item \textbf{Implicit feedback.} A preference inferred from observed behavior.
\item \textbf{Exposure.} Evidence that an item was actually shown and could be acted on.
\end{itemize}
